\clearpage\corrigeDevoir{2}
\begin{blocLivre}{dscor-2-1}\begin{corrigeBac}{6}{Un ion à reconstituer}
\demarche{Identifier le noyau avant de compter le déséquilibre de charges.}
\begin{questions}
\item $A=13+14=27$ ; $\ce{^{27}_{13}Al}$.\pts{2 pt}
\item $q=(13-10)e=+3e=\qty{4,80e-19}{C}$ ; $\ce{Al^{3+}}$.\pts{2 pt}
\item Un ion aluminium pour trois chlorure : $\ce{AlCl3}$.\pts{2 pt}
\end{questions}
\end{corrigeBac}\end{blocLivre}
\begin{blocLivre}{dscor-2-2}\begin{corrigeBac}{7}{Des électrons aux molécules}
\demarche{Distinguer électrons de valence de l’atome isolé et électrons entourant l’atome lié.}
\schema[corrige]{lewis}
\begin{questions}
\item $1s^2\,2s^2\,2p^4$ ; six électrons de valence.\pts{2 pt}
\item Deux doublets liants et deux non liants donnent huit électrons autour de O ; chaque H a un duet.\pts{2 pt}
\item Celle du néon, à couche externe saturée.\pts{1 pt}
\item La liaison stabilise l’assemblage par rapport aux atomes séparés ; sa rupture demande l’énergie de liaison.\pts{2 pt}
\end{questions}
\end{corrigeBac}\end{blocLivre}
\begin{blocLivre}{dscor-2-3}\begin{corrigeBac}{7}{Compter du dioxyde de carbone}
\demarche{Calculer la masse d’une entité puis employer les mêmes unités pour les deux masses.}
\begin{questions}
\item $m_e=m_C+2m_O=\qty{7,31e-26}{kg}$.\pts{2 pt}
\item $N=7,31\times10^{-4}/(7,31\times10^{-26})=10^{22}$ ; $n=\qty{0,0166}{mol}$.\pts{2 pt}
\item $2N=\num{2,00e22}$ atomes d’oxygène.\pts{1 pt}
\item On mélangerait g et kg et on surestimerait N par un facteur 1000.\pts{2 pt}
\end{questions}
\end{corrigeBac}\end{blocLivre}

